% GENERATED FILE. Edit canonical lessons, then run npm run generate:books.
% canonical-source-hash: fnv1a64:2bda5d4ec0a41955
% canonical-lessons: ES-C716-lastima, ES-C716-consolar, ES-C716-lamentar, ES-C716-decepcionado, ES-C716-decepcion

\chapter{Malas noticias}
\label{ch:es-malas-noticias}
\hlchaptermodality{716}

\begin{chapteropening}
\textbf{By the end of this chapter:} \emph{I can react to bad news, show pity and comfort someone who is disappointed.}

Complete the last lesson of chapter 716: I can react to bad news, show pity and comfort someone who is disappointed.
\end{chapteropening}

\section[la lástima]{la lástima — pity, shame}

\label{lesson:ES-C716-lastima}



\begin{quote}
Before the new one: say the Spanish for to kiss, then the Spanish for passion.
\end{quote}

\subsection*{You'll want to know: la lástima}
\textbf{la lástima} — \textquotedblleft{}pity, shame\textquotedblright{}.

\begin{quote}
¡Qué lástima que no vengas!
\end{quote}

\begin{quote}
\emph{What a shame you are not coming!}
\end{quote}

\emph{¡Qué lástima!} — What a pity! \emph{Qué lástima que} takes the subjunctive.

\begin{cousinweb}
\emph{lástima} is built on \emph{lastimar} \textquotedblleft{}to hurt\textquotedblright{}, which comes from a spoken Latin \emph{blastemāre}, Latin \emph{blasphēmāre} \textquotedblleft{}to insult\textquotedblright{}, from Greek \emph{blasphēmía} (EtymDB). The sense moved from insulting to hurting to pity.
\end{cousinweb}

\begin{grammarlens}[title={What you've built}]
One more word for telling, explaining and describing at B1.
\end{grammarlens}

\subsection*{Your turn}
\begin{itemize}
\raggedright
  \item \emph{Say it:} \emph{la lástima}
  \item \emph{Say it:} \emph{la lástima}, once more
  \item \emph{Say it:} a sentence of your own with \emph{la lástima}
  \item \emph{Recall:} say the Spanish for to kiss, then the Spanish for passion, then say \emph{la lástima} again
\end{itemize}

\begin{tcolorbox}[breakable,colback=teal!4,colframe=teal!35!black,title={Before you move on}]
What does la lástima mean? (\textquotedblleft{}Pity, shame\textquotedblright{}.) Say it once more.
\end{tcolorbox}
\section[consolar]{consolar — to console, to comfort}

\label{lesson:ES-C716-consolar}



\begin{quote}
Before the new one: say the Spanish for passion, then the Spanish for pity.
\end{quote}

\subsection*{You'll want to know: consolar}
\textbf{consolar} — \textquotedblleft{}to console, to comfort\textquotedblright{}.

\begin{quote}
Consolé a mi amiga después del examen.
\end{quote}

\begin{quote}
\emph{I comforted my friend after the exam.}
\end{quote}

Stem change: \emph{consuelo}, \emph{consuelas}.

\begin{cousinweb}
\emph{consolar} is inherited from Latin \emph{cōnsōlārī} \textquotedblleft{}to console, to offer comfort\textquotedblright{} (EtymDB).
\end{cousinweb}

\begin{grammarlens}[title={What you've built}]
One more word for telling, explaining and describing at B1.
\end{grammarlens}

\subsection*{Your turn}
\begin{itemize}
\raggedright
  \item \emph{Say it:} \emph{consolar}
  \item \emph{Say it:} \emph{consolar}, once more
  \item \emph{Say it:} a sentence of your own with \emph{consolar}
  \item \emph{Recall:} say the Spanish for passion, then the Spanish for pity, then say \emph{consolar} again
\end{itemize}

\begin{tcolorbox}[breakable,colback=teal!4,colframe=teal!35!black,title={Before you move on}]
What does consolar mean? (\textquotedblleft{}To console, to comfort\textquotedblright{}.) Say it once more.
\end{tcolorbox}
\section[lamentar]{lamentar — to regret, to be sorry about}

\label{lesson:ES-C716-lamentar}



\begin{quote}
Before the new one: say the Spanish for pity, then the Spanish for to console.
\end{quote}

\subsection*{You'll want to know: lamentar}
\textbf{lamentar} — \textquotedblleft{}to regret, to be sorry about\textquotedblright{}.

\begin{quote}
Lamento llegar tarde.
\end{quote}

\begin{quote}
\emph{I am sorry I am late.}
\end{quote}

Formal: \emph{lamentamos las molestias} — we apologise for the inconvenience.

\begin{cousinweb}
\emph{lamentar} comes from Latin \emph{lāmentārī}, built on \emph{lāmentum} \textquotedblleft{}wail, wailing\textquotedblright{} (EtymDB).
\end{cousinweb}

\begin{grammarlens}[title={What you've built}]
One more word for telling, explaining and describing at B1.
\end{grammarlens}

\subsection*{Your turn}
\begin{itemize}
\raggedright
  \item \emph{Say it:} \emph{lamentar}
  \item \emph{Say it:} \emph{lamentar}, once more
  \item \emph{Say it:} a sentence of your own with \emph{lamentar}
  \item \emph{Recall:} say the Spanish for pity, then the Spanish for to console, then say \emph{lamentar} again
\end{itemize}

\begin{tcolorbox}[breakable,colback=teal!4,colframe=teal!35!black,title={Before you move on}]
What does lamentar mean? (\textquotedblleft{}To regret, to be sorry about\textquotedblright{}.) Say it once more.
\end{tcolorbox}
\section[decepcionado]{decepcionado — disappointed}

\label{lesson:ES-C716-decepcionado}



\begin{quote}
Before the new one: say the Spanish for to console, then the Spanish for to regret.
\end{quote}

\subsection*{You'll want to know: decepcionado}
\textbf{decepcionado} — \textquotedblleft{}disappointed\textquotedblright{}.

\begin{quote}
Estoy decepcionado con el hotel.
\end{quote}

\begin{quote}
\emph{I am disappointed with the hotel.}
\end{quote}

A false friend: it means disappointed, not deceived.

\begin{cousinweb}
\emph{decepcionado} is built on \emph{decepción}, a learned borrowing from Latin \emph{dēceptiō} \textquotedblleft{}deceit\textquotedblright{}, from \emph{dēcipere} \textquotedblleft{}to trick\textquotedblright{}; the modern sense of disappointment probably came through French \emph{déception}.
\end{cousinweb}

\begin{grammarlens}[title={What you've built}]
One more word for telling, explaining and describing at B1.
\end{grammarlens}

\subsection*{Your turn}
\begin{itemize}
\raggedright
  \item \emph{Say it:} \emph{decepcionado}
  \item \emph{Say it:} \emph{decepcionado}, once more
  \item \emph{Say it:} a sentence of your own with \emph{decepcionado}
  \item \emph{Recall:} say the Spanish for to console, then the Spanish for to regret, then say \emph{decepcionado} again
\end{itemize}

\begin{tcolorbox}[breakable,colback=teal!4,colframe=teal!35!black,title={Before you move on}]
What does decepcionado mean? (\textquotedblleft{}Disappointed\textquotedblright{}.) Say it once more.
\end{tcolorbox}
\section[la decepción]{la decepción — disappointment}

\label{lesson:ES-C716-decepcion}



\begin{quote}
Before the new one: say the Spanish for to regret, then the Spanish for disappointed.
\end{quote}

\subsection*{You'll want to know: la decepción}
\textbf{la decepción} — \textquotedblleft{}disappointment\textquotedblright{}.

\begin{quote}
La película fue una decepción.
\end{quote}

\begin{quote}
\emph{The film was a disappointment.}
\end{quote}

A false friend: it means disappointment, not deception.

\begin{cousinweb}
\emph{decepción} is a learned borrowing from Latin \emph{dēceptiō} \textquotedblleft{}deceit\textquotedblright{}, from \emph{dēcipere} \textquotedblleft{}to catch out\textquotedblright{}, built on \emph{capere} \textquotedblleft{}to take\textquotedblright{}; it probably took the sense of disappointment from French \emph{déception}.
\end{cousinweb}

\begin{grammarlens}[title={What you've built}]
One more word for telling, explaining and describing at B1.
\end{grammarlens}

\subsection*{Your turn}
\begin{itemize}
\raggedright
  \item \emph{Say it:} \emph{la decepción}
  \item \emph{Say it:} \emph{la decepción}, once more
  \item \emph{Say it:} a sentence of your own with \emph{la decepción}
  \item \emph{Recall:} say the Spanish for to regret, then the Spanish for disappointed, then say \emph{la decepción} again
\end{itemize}

\begin{tcolorbox}[breakable,colback=teal!4,colframe=teal!35!black,title={Before you move on}]
What does la decepción mean? (\textquotedblleft{}Disappointment\textquotedblright{}.) Say it once more.
\end{tcolorbox}
